% !TEX root = ../../proposal.tex

\section{Generative Models for Medical Imaging}
\label{sec:rw_medical}

Deep generative models have been applied to chest radiography synthesis along several
axes. GAN-based approaches~\citep{gan_medical_image_augmentation_2018,
medical_image_synthesis_gans_2018} demonstrated that synthetic images can augment small
annotated datasets and improve classifier performance, but suffer from mode collapse and
limited control over pathology location and type. Diffusion-based models have improved
both fidelity and diversity: Kazerouni et al.~\citep{kazerouni2022diffusion} surveyed
the growing use of diffusion models across medical imaging modalities, and Hashmi et
al.~\citep{hashmi2024xreal} showed high-quality X-ray synthesis using latent diffusion.
These methods condition on a single pathology label or text description; they produce
plausible single-disease images but cannot synthesise a coherent co-morbid presentation
because they have no mechanism for independently specifying and combining two disease
conditions. This is the central limitation Phase~3 addresses.

CheSS~\citep{NEURIPS2022_123d3e81} trained a chest X-ray encoder using a
self-supervised contrastive objective on large radiograph collections, producing
domain-calibrated embeddings. This is directly relevant to Phase~3: unlike general-purpose
CLIP, CheSS similarity scores reflect radiologically meaningful differences between
images, making it appropriate both for conditioning per-disease experts and for evaluating
whether composed images are radiologically consistent with real co-morbid cases. The
VinBigData dataset~\citep{nguyen2022vindr}, used in Phase~3, provides the clinical
setting: 15{,}000 PA radiographs with 14-class radiologist annotations and bounding box
localisation, offering both pathology labels and spatial priors that guide the per-disease
latent factorisation. SepVAE~\citep{louiset2024sepvae}, discussed in
Section~\ref{sec:rw_disentanglement}, is also relevant here: it demonstrated latent
partitioning for pathology separation in a contrastive cohort setting. Phase~3 extends
this idea to a compositional setting in which two co-existing pathologies must be
simultaneously represented and combined, without a contrastive healthy/pathological split.
No prior work has applied score-level composition with factorised disease latents to
co-morbid chest radiograph synthesis; Phase~3 is the first to do so.
